\section{问题重述}

\subsection{问题背景}

山区洪涝灾害具有突发性强、影响范围分散以及道路、供电和通信设施同时受损的特点。连续强降雨可能引发山洪、滑坡和道路塌方，使部分居民点在短时间内与外界失去稳定的地面交通联系。灾后早期，饮用水、应急食品、医疗物资和生活卫生用品需要尽快送达，而传统车辆往往难以及时进入地形起伏大、道路受损严重的区域；无人机受地面道路条件影响较小，可作为应急运输的重要补充手段\cite{ref_xinhua,ref_dianxun}。

2026 年 7 月，受台风"美莎克"带来的持续强降雨影响，广西横州市镇龙乡遭遇严重洪涝灾害，一度出现断路、断电、断网，全乡约 8000 人受困。救援期间，食品、饮用水和药品需依靠徒步或空中投送进入受灾区域。这一现实场景表明：应急物资运输和通信保障需要在道路受阻、能源受限和信息不完整的条件下协同组织。

在完成物资投送的过程中，无人机运输能力受载荷、水平飞行距离、爬升与下降高度、沿线地形净空、能源储备、任务时限和设备周转等多种因素制约；同时，运输无人机在爬升、巡航、下降及物资交接的全过程需持续保持指挥、遥测和交接确认链路。当固定通信设施受损或山体遮挡导致地面网关无法覆盖完整任务区间时，还需空中中继无人机在适当位置和时段提供临时通信保障。因此，山区洪涝灾害下的物资投送既涉及无人机运输，也需同步考虑通信保障和设备资源的协调使用。

\subsection{问题内容}

题目以镇龙乡及周边山区为地理背景，设置标准测试场景：1 个临时调度中心 O01、15 个服务区 S001--S015、80 个不可拆货箱，配置 A/B/C 三种运输无人机机型共 8 架实体机及配套共享电池组，并提供 30\,m DEM 数据、中继无人机、可更换能源组件及无线链路参数等。本文据此研究无人机物资运输、通信中继与救援资源配置之间的协同关系，需依次解决四个递进问题。

\textbf{问题一（单点往返运输能力与货箱组批方案）：} 在不考虑实体无人机和共享电池调度的前提下，每个架次采用 O01$\to$Si$\to$O01 的直接往返形式，仅服务一个服务区。要求：(1) 计算三种机型在不同服务区执行单点往返任务时的最大安全载荷 $q_{\max}$；(2) 在货箱不可拆分、每箱只安排一次，并满足载质量、装载体积和返航安全能量余量的条件下，确定各服务区的货箱组批方案；(3) 在完成全部货箱交付的基础上，综合考虑往返架次数、总运输能耗和累计作业时间，对组批方案进行多目标优化并说明指标间的权衡关系；(4) 讨论返航安全余量变化对最大安全载荷和组批结果的影响。

\textbf{问题二（异构无人机多点多架次运输调度）：} 在不考虑通信保障的条件下，每个运输架次可访问一个或多个服务区，并在完成交付后返回 O01。在满足货箱不可拆分、载质量与装载体积、返航安全余量、无人机与共享电池数量、充电周转及物资时限等约束下，联合确定货箱组批、服务区访问顺序、运输机型、执行无人机、共享电池及各架次开始时刻。其中医疗物资须满足期望送达时间，首批保障货箱须满足首批截止时间（硬时限），其余物资的期望送达时间用于衡量配送及时性。要求：(1) 综合考虑配送及时性、全部任务完成时间、运输能耗和架次数进行优化并说明权衡关系；(2) 给出运输路线、架次安排、逐箱送达时刻及资源使用情况并检验资源可行性。

\textbf{问题三（通信约束下的运输与中继联合调度）：} 在考虑通信保障的条件下，研究运输与中继无人机的联合调度。运输无人机在爬升、巡航、下降及物资投送阶段均须保持连续通信，直连不可用时由空中中继提供保障，通信状态按附录 3 及配套参数判定。要求在满足货箱时限、载荷与能量、连续通信及资源可用性等约束下，联合确定组批与访问顺序、运输及电池安排、架次开始时刻，以及中继无人机的悬停位置、飞行海拔、服务时段和能源组件分配。要求：(1) 综合考虑配送及时性、联合任务完成时间、运输与中继总能耗、两类无人机架次数进行优化并说明权衡；(2) 给出运输与中继联合调度方案并说明货箱交付、资源使用和通信保障情况。

\textbf{问题四（救援任务分区与资源配置优化）：} 以问题三的联合调度方案为基础，将 15 个服务区划分为 2 个和 3 个任务组。每个服务区必须且只能属于一个任务组；分区过程中保持问题三确定的组批、访问顺序、运输与中继任务安排及通信保障关系不变；同一运输架次涉及多个服务区时须划入同一任务组；各组仅承担本组服务区对应的任务并分别核算独立执行所需的运输无人机、共享电池、中继无人机和中继能源组件数量，资源不得跨组调配。要求：(1) 给出 K=2/K=3 的分区及资源配置方案，从资源配置规模、资源冗余、组间工作量均衡和资源缺口等方面比较两种分区方式；(2) 对资源缺口给出原因分析。

\subsection{各问题的数据继承关系}

本文四个问题构成逐问递进的求解链条：问题一建立公共的"地形净空—能耗—时间"引擎并完成单点组批；问题二在问题一基础上引入实体机队、共享电池池与多站路线，形成运输调度；问题三在问题二运输方案上叠加通信中继维度，形成运输与中继联合调度；问题四以问题三方案为基底，只做任务分区与资源核算、不改动既有调度关系。全文共用同一套 DEM 采样、能耗公式、两阶段充电与作业时间口径，保证了各问结果的可比性与可继承性。
